% 原书补充：SGA/SGA-7-2-Groupes-de-monodromie-2.pdf，PDF 第 5--7 页。
% 此引言未包含于所选英文 release TeX；由法文原页逐段译出。
\clearpage
\begin{center}
\underline{引言}
\end{center}
\medskip

本讨论班由 A. Grothendieck 主持的第一部分，已收入同一套 Lecture Notes in Mathematics 的第 288 卷。对于第一部分，我们只需要使用第 I 讲和第 VI 讲（特别是 \S\S5、6）的结果。

这第二部分的关键结果，是第 XV 讲的 Picard--Lefschetz 公式，以及第 XVIII 讲的 Lefschetz 线束理论。

设 $X_t$ 是依赖于参数 $t$ 的一族代数簇。Picard--Lefschetz 公式描述 $X_t$ 的上同调在参数的某个值 $t_0$ 附近的行为；在该值处，$X_t$ 出现一个普通二次奇点：该公式给出 $X_t$ 与 $X_{t_0}$ 的上同调之间的差异，并描述 $t$ 绕 $t_0$ 转动时的局部单值。在超越情形，此公式由 Lefschetz [2] 得到，具有简单的几何意义（XIV 3.2）。为在任意特征下建立它，我们必须构筑一套庞大的机制（XII、XIII、XIV）；其轮廓见第 XIII、XIV 讲的引言。

当取非奇异射影簇 $X$ 的足够一般的超平面截面线束作为族 $X_t$ 时，Picard--Lefschetz 公式连同其他几何思想，给出 $X$ 的上同调与其各超平面截面 $X_t$ 的上同调之间的宝贵联系。只有当 $X$ 满足强 Lefschetz 定理（LV）（XVIII 5.2.2）时，这些联系才能发挥全部效力。特征为 $0$ 时，这一条件自动成立。当维数 $\geq3$ 时，唯一已知的证明依赖于调和积分的 Hodge 理论（Lefschetz 对关键之处的证明 [2] 13, p.~25 有误）；若能找到另一个可推广至特征 $p$ 的证明，将极其有意义。

主要应用如下。

\begin{enumerate}[label=\alph*),leftmargin=2.6em]
\item 在第 XVI 讲中，我们在等特征 $p$ 下，证明了一个类似于 Milnor [3] 对映射 $f:\mathbb C^n\longrightarrow\mathbb C$ 的孤立临界点处的消灭圈数目所给出的公式。

\item 在第 XIX 讲中，我们使用 Lefschetz 的方法（[2] 第 106 页），特别是在任意特征下证明：若 $S$ 是 $\mathbb P^3$ 中次数 $\ne2,3$ 的泛曲面，则 $S$ 上的每条曲线都是 $S$ 与 $\mathbb P^3$ 中另一个曲面的完全交。

\item 在第 XX 讲中，我们将 Griffiths [1] 的方法转用到任意特征下，证明存在同调等价于零的代数圈，而其任何非零倍数都不代数等价于零。
\end{enumerate}

还要指出，Lefschetz 线束方法最近使人们得以按维数归纳证明如下结果（见 [4]）。若 $\mathbb F_p$ 上的非奇异射影簇 $X$ 能（连同其极化）提升至特征 $0$，且 $p\ne2$，则 Frobenius 作用于 $H^i(X\otimes_{\mathbb F_p}\overline{\mathbb F}_p,\mathbb Q_\ell)$ 的特征多项式属于 $\mathbb Z[t]$，并且不依赖于 $\ell$（$\ell\ne p$）。

第 XXII、XXI 讲是第 XX 讲的预备知识（在特征 $2$ 时，也是第 XIX 讲的预备知识）。在第 XXII 讲中，我们研究 $\mathbb F_p$ 上代数簇的 $\zeta$ 函数的模 $p$ 约化（所用方法中 Dwork 理论起关键作用），以及 Frobenius 作用于 $H^i$ 的特征多项式的整性性质。在第 XXI 讲中，我们给出泛完全交的上同调的层次的下界。
% 源文 PDF 第 6 页此处印作 propriétés d'intégrabilité；译文按本句
% 及第 XXII 讲的数学对象理解为特征多项式的整性，详见 review/frontmatter.json。

对于希望先获得这一理论的几何印象的读者，我建议首先阅读第 I 讲、XIV 3.2、XIX \S4，以及第 XIII 讲的引言。

后面的 Leitfaden 大致说明了各讲之间的逻辑依赖关系。

\medskip
\begin{enumerate}[label={[\arabic*]},leftmargin=2.8em]
\item P.A. GRIFFITHS，\emph{On the periods of certain integrals III}，Ann. of Math.
\item S. LEFSCHETZ，\emph{L'analyse situs et la g\'eom\'etrie alg\'ebrique}，Paris 1924，Gauthier-Villars，重印于 \emph{Selected Papers}（Chelsea）。
\item J. MILNOR，\emph{Singular points of complex hypersurfaces}，Annals of Math. Studies 61，Princeton 1968。
\item J.L. VERDIER（据 P. DELIGNE），\emph{Cohomologie $\ell$-adique}，S\'eminaire Bourbaki，1972 年 11 月。
\end{enumerate}

\vspace{2em}
\begin{flushright}
Bures-sur-Yvette，1972 年 11 月 8 日\\[1.5em]
P. DELIGNE
\end{flushright}
